\renewcommand{\bibname}{Quellen und Bandzuordnung}
\begin{thebibliography}{9}
\addcontentsline{toc}{chapter}{Quellen und Bandzuordnung}
\footnotesize
\setlength{\itemsep}{4pt}
\setlength{\parskip}{0pt}
\interlinepenalty=10000

\bibitem{project}
Martin Kunze: \emph{Die Grundlagen der Mathematik}.
Band~7: Begriff der Surjektivität, volles Bild, Cantors Satz,
Projektionen, Fasern, Einschränkungen, Komposition und eindeutige
Faktorisierung. Band~3 und Band~5 liefern die Mengen- und
Funktionsbegriffe. Format nach den Leitfäden zu Band~3,
Band~4, Band~5, Band~28 und Band~37.

\bibitem{hammack}
Richard Hammack: \emph{Book of Proof}.
3.~Auflage, Fassung~3.4.
Abschnitt~12.2, Definition~12.4, Druckseite~228:
Surjektivität. Abschnitt~12.4, Satz~12.2, Druckseite~237:
Komposition surjektiver Funktionen.
Abschnitt~12.6, Definition~12.9 und Aufgabe~12,
Druckseiten~242--243: Bild, Urbild und das Kriterium
$f[f^{-1}[V]]=V$ für alle Zielteilmengen.
\url{https://richardhammack.github.io/BookOfProof/Main.pdf}.

\bibitem{eisermann}
Michael Eisermann: \emph{Lineare Algebra}, Kapitel~H,
Halbversion vom 18.~September 2026, Universität Stuttgart.
Satz~H3l, Folie~H344: Faktorisierung über eine Surjektion,
Eindeutigkeit und Gleichheit der Bilder.
Verwendet wird ausschließlich die mengentheoretische Aussage,
hier unmittelbar durch Konstanz auf den Fasern formuliert.
Die übrigen Themen dieses Kapitels werden nicht übernommen.
\url{https://pnp.mathematik.uni-stuttgart.de/igt/eiserm/lehre/Lineare-Algebra/LinA-halb-H-2x4.pdf}.

\bibitem{moller}
Jesper M.~Møller: \emph{General Topology},
Universität Kopenhagen, Kapitel~1: \emph{Sets and maps}.
Satz~1.28, Druckseite~14: Zusammenhang des Auswahlaxioms
mit Rechtsinversen von Surjektionen und mit Produkten
nichtleerer Mengen. Diese Quelle begründet allein die
Abgrenzung gegenüber Auswahlannahmen; topologische Strukturen
und Auswahlprinzipien sind keine Gegenstände dieses Leitfadens.
\url{https://web.math.ku.dk/~moller/e03/3gt/notes/gtnotes.pdf}.

\end{thebibliography}

\noindent\textbf{Einordnung der Zusammenstellung.}
Die Quellen dienen der Kontrolle der Begriffe und Kernaussagen.
Die konkreten Zusammenfassungen, die Behandlung leerer Mengen
und die ergänzenden Konstruktionen sind eigenständig formuliert.
Nicht jede Aussage steht in den Quellen in genau dieser Gestalt.
Der Leitfaden beansprucht weder historische Erstentdeckungen
noch eine vollständige Erfassung aller bekannten Resultate.
Er enthält ausschließlich Definitionen, Sätze und kurze Hinweise
zu ihrem Geltungsbereich.